\documentclass[11pt, a4paper]{article}

\usepackage[T1]{fontenc}
\usepackage[utf8]{inputenc}
\usepackage{tgpagella}
\usepackage[spanish, es-noshorthands]{babel}
\usepackage{amsmath, amssymb, bm}
\usepackage{tabularx}
\usepackage{booktabs}
\usepackage[most]{tcolorbox}
\usepackage[margin=2.5cm]{geometry}
\usepackage{ragged2e}
\usepackage{fancyhdr}
\usepackage[american]{circuitikz}

\pagestyle{fancy}
\fancyhf{}
\setlength{\headheight}{16pt}
\lhead{\textbf{UCB Sede Tarija} | Ing. Mecatrónica}
\chead{}
\rhead{IMT-121 Circuitos Electrónicos I | S08-02}
\lfoot{Prof: M.Sc. Ing. Bernardo Quiroga Turdera}
\rfoot{Página \thepage}
\renewcommand{\headrulewidth}{0.4pt}

\definecolor{ucbblue}{RGB}{0, 51, 102}

\begin{document}

\begin{tcolorbox}[colback=ucbblue!5!white, colframe=ucbblue, title=\textbf{Consolidación Práctica Previa al Examen (Hito 2)}]
    \centering \textbf{Superposición, Thévenin, Norton y Máxima Transferencia}
\end{tcolorbox}

\noindent \textit{Instrucciones: Trabaje primero de forma individual. Justifique sus referencias, la desactivación de fuentes y su método. La respuesta numérica sin procedimiento no constituye evidencia de comprensión. Los cálculos están diseñados para resolverse sin calculadora.}

\vspace{0.4cm}
\section*{Problema 1: Superposición y Referencias}
Dado el siguiente circuito lineal, se desea conocer el voltaje $V_a$ utilizando el principio de superposición.

\begin{center}
\begin{circuitikz}[scale=1, transform shape]
    \draw (0,0) to[V, l=$12\,\text{V}$, invert] (0,2) to[R, l=$2\,\text{k}\Omega$] (3,2);
    \draw (6,0) to[V, l_=$6\,\text{V}$, invert] (6,2) to[R, l_=$6\,\text{k}\Omega$] (3,2);
    \draw (3,2) to[short, *-o] (3,3) node[above]{$a$};
    \draw (3,2) to[R, l=$3\,\text{k}\Omega$] (3,0);
    \draw (0,0) -- (6,0);
    \draw (3,0) node[ground]{};
\end{circuitikz}
\end{center}

\begin{enumerate}
    \item Dibuje el circuito equivalente dejando activa \textbf{solamente} la fuente de $12\,\text{V}$ y calcule la contribución $V_a^{(1)}$.
    \item Dibuje el circuito equivalente dejando activa \textbf{solamente} la fuente de $6\,\text{V}$ y calcule la contribución $V_a^{(2)}$.
    \item Obtenga el voltaje total $V_a$. Indique explícitamente la referencia respecto a la cual está medido.
    \item \textbf{Análisis:} ¿Qué significaría físicamente si una de las contribuciones resultara tener signo negativo?
\end{enumerate}

\section*{Problema 2: Equivalencia de Puerto (Thévenin y Norton)}
Se tiene una red de alimentación conectada a una resistencia de carga $R_L$ en el puerto $a-b$ (siendo $b$ el nodo de tierra). La red interna consiste en una fuente ideal de $12\,\text{V}$ en serie con una resistencia $R_1 = 4\,\text{k}\Omega$ conectada al nodo $a$, y una resistencia $R_2 = 4\,\text{k}\Omega$ en paralelo entre el nodo $a$ y tierra.

\begin{enumerate}
    \item Retire conceptualmente la carga $R_L$ y calcule el voltaje de circuito abierto ($V_{th}$).
    \item Redibuje la red con las fuentes desactivadas correctamente y determine $R_{th}$ mirando desde el puerto $a-b$.
    \item Dibuje el equivalente de Thévenin resultante.
    \item Calcule $I_N$ y dibuje el equivalente de Norton.
\end{enumerate}

\section*{Problema 3: Máxima Transferencia y Efecto de Carga}
Considerando el equivalente de Thévenin obtenido en el Problema 2 ($V_{th} = 6\,\text{V}$, $R_{th} = 2\,\text{k}\Omega$), evalúe el comportamiento del circuito al conectar sucesivamente las siguientes cargas: $1\,\text{k}\Omega$, $2\,\text{k}\Omega$ y $4\,\text{k}\Omega$.
\begin{enumerate}
    \item Calcule $I_L$, $V_L$ y $P_L$ para cada una de las tres cargas.
    \item Responda sin realizar cálculos adicionales:
    \begin{itemize}
        \item ¿Qué carga produce la mayor corriente?
        \item ¿Qué carga produce el mayor voltaje en terminales?
        \item ¿Cuál produce la mayor potencia útil? ¿Por qué no coinciden las tres respuestas?
    \end{itemize}
\end{enumerate}

\section*{Problema 4: Interpretación de Datos Experimentales}
Un equipo reporta la siguiente tabla de mediciones (datos simulados) obtenida de una red cuyo modelo teórico indica $R_{th} = 4\,\text{k}\Omega$.

\vspace{0.2cm}
\noindent
\begin{tabularx}{\textwidth}{*{4}{>{\centering\arraybackslash}X}}
\toprule
\textbf{$R_L$ (k$\Omega$)} & \textbf{$V_L$ medido (V)} & \textbf{$I_L$ medido (mA)} & \textbf{$P_L$ medido (mW)} \\
\midrule
$2$ & $3.9$ & $1.95$ & $7.60$ \\ \addlinespace
$3$ & $5.0$ & $1.67$ & $8.35$ \\ \addlinespace
$4$ & $5.8$ & $1.45$ & $8.41$ \\ \addlinespace
$5$ & $6.5$ & $1.30$ & $8.45$ \\ \addlinespace
$8$ & $7.8$ & $0.98$ & $7.64$ \\
\bottomrule
\end{tabularx}
\vspace{0.2cm}

\begin{enumerate}
    \item Según los datos de la tabla, ¿en qué rango se encuentra el máximo experimental de potencia?
    \item ¿Cómo se compara este hallazgo con el $R_{th}$ teórico de $4\,\text{k}\Omega$?
    \item ¿Esta ligera desviación invalida la teoría analítica? Proponga \textbf{una hipótesis técnica plausible} que justifique la diferencia observada y sugiera qué medición adicional le permitiría comprobar dicha hipótesis.
\end{enumerate}

\end{document}
